% ============================================
% Compatible with pdfLaTeX and XeLaTeX
% ============================================

\documentclass[12pt,a4paper]{article}

% ---------- Encoding and Language ----------
\usepackage[utf8]{inputenc}          % UTF-8 encoding (use utf8x if extended Cyrillic needed)
\usepackage[T2A]{fontenc}           % Cyrillic font encoding (for pdfLaTeX)
\usepackage[russian]{babel}         % Russian language support

% ---------- Page Geometry ----------
\usepackage[
  left=2cm,
  right=2cm,
  top=2cm,
  bottom=2cm,
  headsep=0.5cm,
  footskip=1cm
]{geometry}

% ---------- Mathematics ----------
\usepackage{amsmath}                % Core math features
\usepackage{amssymb}                % Additional math symbols
\usepackage{amsthm}                 % Theorem-like environments

% ---------- Lists (Compact & Multi-level) ----------
\usepackage{enumitem}
% Set compact defaults for all lists
\setlist{
  noitemsep,                        % No extra space between items
  topsep=0.5ex,                     % Minimal space above/below list
  partopsep=0pt,                    % No extra space when starting a paragraph
  parsep=0pt,                       % No space between paragraphs within items
  leftmargin=2em,                   % Consistent indentation
  labelsep=0.5em                    % Space between label and text
}
% Fine-tune specific list types
\setlist[itemize]{label=---}        % Em-dash bullets
\setlist[enumerate,1]{label=\arabic*.} % First level: 1.
\setlist[enumerate,2]{label=\alph*)}  % Second level: a)
\setlist[enumerate,3]{label=\roman*.} % Third level: i.
\setlist[description]{font=\normalfont\itshape} % Description labels in italics

% ---------- Problem Environment ----------
\newcounter{task}                   % Custom counter for problems
\newcommand{\taskname}{}            % Placeholder for task label

% Define the task environment
\newenvironment{task}[1][]{%
  \refstepcounter{task}%            % Increment and reference the counter
  \par\vspace{0.3\baselineskip}%    % Tiny space before problem
  \noindent\textbf{\thetask.}%      % Bold number with period
  \if\relax\detokenize{#1}\relax   % Check if optional title is provided
    \ \ignorespaces                 % Just space after number
  \else
    \ \textbf{#1}.\ \ignorespaces   % Space, bold title, period, space
  \fi
}{%
  \par\vspace{0.2\baselineskip}%    % Tiny space after problem
}

% ---------- Solution Command ----------
\newcommand{\sol}{%
  \par\vspace{0.3\baselineskip}%    % Small space before solution
  \noindent\textbf{Решение.}\ %    % Bold "Решение." with space after
  \ignorespaces                     % Ignore spaces after the command
}

% ---------- Solution Sub-item Command (for numbered parts) ----------
% Usage: \soli{1} or \soli{a}
\newcommand{\soli}[1]{%
  \par\vspace{0.15\baselineskip}%   % Tiny space before sub-solution
  \noindent\textbf{#1)}\ %          % Bold number/letter with parenthesis
  \ignorespaces                     % Ignore spaces after the command
}

% ---------- Header Configuration ----------
\makeatletter
\renewcommand{\@maketitle}{%
  % Top row: Sheet number (left) and Student name (right)
  \begin{flushleft}
    \textbf{Листок \#\sheetnumber}%
    \hfill%
    \textbf{\studentname}%
  \end{flushleft}
  \vspace{-0.8\baselineskip}%      % Tighten vertical spacing
  
  % Bottom row: Assignment title
  \begin{flushleft}
    \textit{\assignmenttitle}%
  \end{flushleft}
  \vspace{-0.5\baselineskip}%      % Additional tightening
}
\makeatother

% Define header variables
\newcommand{\sheetnumber}[1]{\def\sheetnumber{#1}}
\newcommand{\studentname}[1]{\def\studentname{#1}}
\newcommand{\assignmenttitle}[1]{\def\assignmenttitle{#1}}

% ---------- Other Settings ----------
\date{}                             % Suppress automatic date
\pagestyle{plain}                   % Simple page numbering
\setlength{\parindent}{0pt}         % No paragraph indentation (modern style)
\setlength{\parskip}{0.5ex}         % Small paragraph spacing

% ============================================
% DOCUMENT BEGIN
% ============================================

\begin{document}

% ---------- Custom Header ----------
% Set your values here:
\sheetnumber{2}                     % Replace with actual sheet number
\studentname{Амехин Илья}           % Replace with student name
\assignmenttitle{Многочлены и конечные поля} % Replace with assignment title
\maketitle


% --- Problem 1 ---
\begin{task}
    Найдите обратный к $f(x)$ в кольце вычетов $k[x]/(g(x))$, если такой существует
    \begin{enumerate}
	\item $\textbf{k} = \mathbb{Q}$, $f(x) = x^6 + x + 1, g(x) = x^2 + 1$
	\item $\textbf{k} = \mathbb{F}_2$, $f(x) = x^7 + 1, g(x) = x^2 + x + 1$
    \end{enumerate}
\end{task}

\sol
\soli{1} Найдем НОД$(f,g)$ в $\mathbb{Q}[x]$. Так как $g(x)=x^2+1$, проверим значения в корнях $g$: $i$ и $-i$. 
\[
f(i)=i^6+i+1=-1+i+1=i\neq 0,\quad f(-i)=(-i)^6-i+1=-1-i+1=-i\neq 0,
\]
$g(x) = x^2 + 1$ неприводим над $\mathbb{Q}$, поэтому из $f(i) \neq 0$ и $f(-i) \neq 0$ следет $\gcd(f(x),g(x)) = 1$
Применим алгоритм Евклида:
\[
\begin{aligned}
f &= x^4 g + (-x^4+x+1),\\
-x^4+x+1 &= (-x^2)g + (x^2+x+1),\\
x^2+x+1 &= 1\cdot g + x,\\
g &= x\cdot x + 1.
\end{aligned}
\]
Отсюда выражаем $1$:
\[
1 = g - x\cdot x = g - x\big((x^2+x+1)-g\big) = (1+x)g - x(x^2+x+1).
\]
Далее, $x^2+x+1 = (-x^4+x+1) - (-x^2)g = (f - x^4 g) + x^2 g = f + (x^2 - x^4)g$. Подставляя, получаем
\[
1 = (1+x)g - x\big(f + (x^2 - x^4)g\big) = -x f + \big(1+x - x(x^2 - x^4)\big)g.
\]
Упрощаем коэффициент при $g$: $1+x - x^3 + x^5 = x^5 - x^3 + x + 1$. Следовательно,
\[
1 = -x\cdot f(x) + (x^5 - x^3 + x + 1)\cdot g(x).
\]
Значит, обратным к $f$ по модулю $g$ является многочлен $-x$. 
\[
\boxed{f^{-1}(x) \equiv -x \pmod{x^2+1}}.
\]

\soli{2} В $\mathbb{F}_2[x]$ имеем $g(x)=x^2+x+1$. Заметим, что $x^3 \equiv 1 \pmod{g(x)}$, так как $x^3+1 = (x+1)(x^2+x+1)$ в $\mathbb{F}_2[x]$. Тогда
\[
x^7 = x^{6}\cdot x = (x^3)^2 \cdot x \equiv 1^2 \cdot x = x \pmod{g},
\]
поэтому
\[
f(x)=x^7+1 \equiv x+1 \pmod{g}.
\]
Остаток $x+1$ не равен нулю и взаимно прост с $g$ (проверка: $g(1)=1$), значит $\gcd(f,g)=1$, обратный существует.

Из $f \equiv x+1$ получаем
\[
x f \equiv x(x+1) = x^2+x \equiv 1 \pmod{g},
\]
поскольку $g = x^2+x+1 \equiv 0 \Rightarrow x^2+x \equiv 1$ (в $\mathbb{F}_2$ знак минус совпадает с плюсом). Следовательно,
\[
\boxed{f^{-1}(x) \equiv x \pmod{x^2+x+1}}.
\]

\newpage 

\begin{task}
    \begin{enumerate}
	\item Пусть $\mathbb{F}_4 \cong \mathbb{F}_2[a]/(\alpha^2 + \alpha + 1)$ -- поле из четырёх элементов, которые мы обозначим $\{0, 1, [\alpha], [\alpha + 1]\}$. Найдите все делители нуля и все обратимые элементы (и обратные к ним) в кольце $\mathbb{F}_4[x] / (x^2 + [\alpha + 1]x + [\alpha])$
    \end{enumerate}
\end{task}

\sol
\soli{1}

Обозначим $a=[\alpha]$. Тогда
\[
a^2+a+1=0,\qquad a^2=a+1,
\]
и
\[
g(x)=x^2+(a+1)x+a.
\]
Имеем
\[
g(1)=1+(a+1)+a=0,
\]
\[
g(a)=a^2+(a+1)a+a=a^2+a^2+a+a=0.
\]
Следовательно,
\[
g(x)=(x+1)(x+a).
\]
Кроме того,
\[
\gcd(x+1,x+a)=1,
\]
поскольку их разность равна $a+1\ne0$. Поэтому по китайской теореме об остатках
\[
\mathbb{F}_4[x]/(g)
\cong
\mathbb{F}_4[x]/(x+1)\times\mathbb{F}_4[x]/(x+a)
\cong\mathbb{F}_4\times\mathbb{F}_4,
\]
где изоморфизм имеет вид
\[
f(x)\longmapsto(f(1),f(a)).
\]

Так как $\deg g=2$, каждый класс имеет единственного представителя
\[
ux+v,\qquad u,v\in\mathbb{F}_4,
\]
и при указанном изоморфизме
\[
ux+v\longmapsto(u+v,\;ua+v).
\]

В произведении $\mathbb{F}_4\times\mathbb{F}_4$ элемент обратим тогда и только тогда, когда обе его компоненты ненулевые, а ненулевой делитель нуля -- тогда и только тогда, когда ровно одна компонента равна нулю.

Если $u+v=0$, то $v=u$, и при $u\ne0$ получаем
\[
x+1,\qquad a(x+1),\qquad (a+1)(x+1).
\]
Если $ua+v=0$, то $v=ua$, и при $u\ne0$ получаем
\[
x+a,\qquad ax+a+1,\qquad (a+1)x+1.
\]
Таким образом, все делители нуля:
\[
\boxed{
x+1,\quad a(x+1),\quad (a+1)(x+1),\quad
x+a,\quad ax+a+1,\quad (a+1)x+1
}.
\]

Обратимых элементов $3\cdot3=9$. Их можно получить из условий
\[
u+v\ne0,\qquad ua+v\ne0.
\]
Получаем:
\[
\begin{array}{c|c}
\text{элемент} & \text{обратный} \\ \hline
1 & 1 \\
a & a+1 \\
a+1 & a \\
x & (a+1)x+a \\
x+a+1 & (a+1)x \\
ax+1 & ax \\
ax & ax+1 \\
(a+1)x & x+a+1 \\
(a+1)x+a & x
\end{array}
\]

Наконец,
\[
1+6+9=16=|\mathbb{F}_4[x]/(g)|,
\]
поэтому перечислены все элементы кольца.

\end{document}
